\documentclass[11pt, a4paper]{article}

% --- UNIVERSAL PREAMBLE BLOCK ---
\usepackage[a4paper, top=2.5cm, bottom=2.5cm, left=2cm, right=2cm]{geometry}
\usepackage{fontspec}

\usepackage[portuguese, bidi=basic, provide=*]{babel}

\babelprovide[import, onchar=ids fonts]{portuguese}
\babelprovide[import, onchar=ids fonts]{english}

% Set default/Latin font to Sans Serif in the main (rm) slot
\babelfont{rm}{Noto Sans}
\babelfont[portuguese]{rm}{Noto Sans}

% Add because main language is not English
\usepackage{enumitem}
\setlist[itemize]{label=-}
% --------------------------------

\usepackage{booktabs}
\usepackage{xcolor}
\usepackage{listings}

% Configuração para exibição do código SQL
\definecolor{codegreen}{rgb}{0,0.6,0}
\definecolor{codegray}{rgb}{0.5,0.5,0.5}
\definecolor{codepurple}{rgb}{0.58,0,0.82}
\definecolor{backcolour}{rgb}{0.95,0.95,0.92}

\lstdefinestyle{mystyle}{
    backgroundcolor=\color{backcolour},   
    commentstyle=\color{codegreen},
    keywordstyle=\color{blue}\bfseries,
    numberstyle=\tiny\color{codegray},
    stringstyle=\color{codepurple},
    basicstyle=\ttfamily\small,
    breakatwhitespace=false,         
    breaklines=true,                 
    captionpos=b,                    
    keepspaces=true,                 
    numbers=none,                    
    showspaces=false,                
    showstringspaces=false,
    showtabs=false,                  
    tabsize=2
}
\lstset{style=mystyle, language=SQL}

\title{\textbf{Relatório Completo de Bases de Dados: Etapas 4 a 7}}
\author{Trabalho Final de Bases de Dados}
\date{}

\begin{document}

\maketitle

Este documento apresenta as etapas finais do projeto prático de Bases de Dados para o sistema da Agência de Turismo, abrangendo Consultas (DML), Atualizações, Criação de Visões e Simulação de Transações. Os resultados apresentados refletem o estado do banco de dados após a execução dos scripts de atualização e normalização (Etapa 5).

% ==============================================================================
% ETAPA 4
% ==============================================================================
\section{Etapa 4: Consultas (SQL/DML)}

\subsection{4(a) Consultas com Junções de Tabelas (JOIN)}

\textbf{Consulta 1: Listar os hotéis, as suas cidades e estados}
\begin{lstlisting}
SELECT 
    H.nomehot AS "Nome do Hotel", 
    C.nomecid AS "Cidade", 
    C.estado AS "Estado"
FROM HOTEL H
JOIN CIDADE C ON H.idcid = C.idcid;
\end{lstlisting}

\begin{table}[h]
\centering
\begin{tabular}{lll}
\toprule
\textbf{Nome do Hotel} & \textbf{Cidade} & \textbf{Estado} \\
\midrule
Hotel Fasano BH & Belo Horizonte & MG \\
Copacabana Palace & Rio de Janeiro & RJ \\
Augusta Inn & São Paulo & SP \\
\bottomrule
\end{tabular}
\end{table}

\textbf{Consulta 2: Mostrar os museus, fundadores e suas profissões}
\begin{lstlisting}
SELECT
    PT.descricaotur AS "Nome do Museu",
    F.nome AS "Fundador",
    F.atividadepro AS "Profissão"
FROM PONTOTURISTICO PT
JOIN MUSEU M ON PT.idtur = M.idtur
JOIN MUSFUND MF ON M.idtur = MF.idtur
JOIN FUNDADOR F ON MF.idpes = F.idpes;
\end{lstlisting}

\begin{table}[h]
\centering
\begin{tabular}{lll}
\toprule
\textbf{Nome do Museu} & \textbf{Fundador} & \textbf{Profissão} \\
\midrule
Museu do Amanhã & Santiago Calatrava & Arquiteto \\
Museu do Amanhã & Eike Batista & Empresário \\
MM Gerdau - Museu das Minas e do Metal & Jorge Gerdau & Empresário \\
\bottomrule
\end{tabular}
\end{table}

\textbf{Consulta 3: Listar os restaurantes, culinária e o hotel onde se localizam}
\begin{lstlisting}
SELECT
    R.nomeres AS "Restaurante",
    R.tipoculi AS "Culinária",
    H.nomehot AS "Hotel"
FROM RESTAURANTE R
JOIN RESTAURANTEHOT RH ON R.idres = RH.idres
JOIN HOTEL H ON RH.idhot = H.idhot;
\end{lstlisting}

\begin{table}[h]
\centering
\begin{tabular}{lll}
\toprule
\textbf{Restaurante} & \textbf{Culinária} & \textbf{Hotel} \\
\midrule
Mee Copacabana Palace & FRA & Copacabana Palace \\
Sushi do Fasano & JAP & Hotel Fasano BH \\
\bottomrule
\end{tabular}
\end{table}

\subsection{4(b) Consultas com Funções de Agregação}

\textbf{Consulta 4: População total das cidades cadastradas}
\begin{lstlisting}
SELECT SUM(populacao) AS "População Total" FROM CIDADE;
\end{lstlisting}
\textit{Resultado:} 21.500.000 habitantes.

\textbf{Consulta 5: Refeição mais cara e mais barata entre os restaurantes}
\begin{lstlisting}
SELECT MAX(precomedioref) AS "Cara", MIN(precomedioref) AS "Barata" FROM RESTAURANTE;
\end{lstlisting}

\begin{table}[h]
\centering
\begin{tabular}{cc}
\toprule
\textbf{Refeição Mais Cara (R\$)} & \textbf{Refeição Mais Barata (R\$)} \\
\midrule
450.00 & 80.00 \\
\bottomrule
\end{tabular}
\end{table}

\textbf{Consulta 6: Média de preço geral das diárias (TIPOQUARTO)}
\begin{lstlisting}
SELECT ROUND(AVG(preco), 2) AS "Média Geral" FROM TIPOQUARTO;
\end{lstlisting}
\textit{Resultado:} R\$ 871.43 (Média calculada sobre todos os tipos disponíveis).

\subsection{4(c) Consultas com Agregação, GROUP BY e HAVING}

\textbf{Consulta 7: Quantidade de quartos por hotel}
\begin{lstlisting}
SELECT H.nomehot AS "Hotel", COUNT(Q.numquarto) AS "Qtd"
FROM HOTEL H JOIN QUARTO Q ON H.idhot = Q.idhot
GROUP BY H.nomehot HAVING COUNT(Q.numquarto) > 2;
\end{lstlisting}

\begin{table}[h]
\centering
\begin{tabular}{lc}
\toprule
\textbf{Hotel} & \textbf{Quantidade de Quartos} \\
\midrule
Hotel Fasano BH & 4 \\
Augusta Inn & 3 \\
Copacabana Palace & 3 \\
\bottomrule
\end{tabular}
\end{table}

\textbf{Consulta 8: Média de diária por hotel (Apenas médias acima de R\$ 400)}
\begin{lstlisting}
SELECT H.nomehot AS "Hotel", ROUND(AVG(TQ.preco), 2) AS "Média"
FROM HOTEL H JOIN TIPOQUARTO TQ ON H.idhot = TQ.idhot
GROUP BY H.nomehot HAVING AVG(TQ.preco) > 400.00;
\end{lstlisting}

\begin{table}[h]
\centering
\begin{tabular}{lc}
\toprule
\textbf{Hotel} & \textbf{Média de Diária (R\$)} \\
\midrule
Hotel Fasano BH & 500.00 \\
Copacabana Palace & 1700.00 \\
\bottomrule
\end{tabular}
\end{table}

\textbf{Consulta 9: Total de pontos turísticos por cidade (Mínimo 2)}
\begin{lstlisting}
SELECT C.nomecid AS "Cidade", COUNT(PT.idtur) AS "Total"
FROM CIDADE C JOIN PONTOTURISTICO PT ON C.idcid = PT.idcid
GROUP BY C.nomecid HAVING COUNT(PT.idtur) >= 2;
\end{lstlisting}

\begin{table}[h]
\centering
\begin{tabular}{lc}
\toprule
\textbf{Cidade} & \textbf{Total de Pontos Turísticos} \\
\midrule
Belo Horizonte & 2 \\
Rio de Janeiro & 2 \\
\bottomrule
\end{tabular}
\end{table}

\clearpage

% ==============================================================================
% ETAPA 5
% ==============================================================================
\section{Etapa 5: Atualizações (SQL/DML)}

Esta seção documenta a implementação de integridade referencial e atualizações complexas.

\subsection{5(a) Inserção com Integridade Referencial}
\begin{lstlisting}
INSERT INTO PONTOTURISTICO (idtur, idcid, tipotur, endereco, descricaotur)
VALUES (6, 1, 'M', 'Rua da Bahia, 1149', 'Museu de Arte da Pampulha');
\end{lstlisting}

\subsection{5(b) Inserção com Cláusula CHECK}
O banco impede a inserção de hotéis com estrelas fora do intervalo [1,5].

\subsection{5(c) Inserção via SELECT (Tabela de Alto Padrão)}
\begin{lstlisting}
INSERT INTO HOTEIS_ALTO_PADRAO (idhot, nomehot, estrelas)
SELECT idhot, nomehot, estrelas FROM HOTEL WHERE estrelas >= 4;
\end{lstlisting}

\subsection{5(d) Remoção com CASCADE}
A remoção do Hotel 3 (\textit{Augusta Inn}) eliminou automaticamente todos os seus quartos e tipos de quarto.

\subsection{5(e) Remoção com SET NULL}
Configurou-se a FK de Cidade em Ponto Turístico para \texttt{SET NULL}. Ao remover o Rio de Janeiro (idcid=2), os pontos turísticos associados permaneceram no banco com a cidade nula.

\subsection{5(f) Atualização por Subconsulta}
Aumento de 15\% nas diárias dos hotéis localizados em Belo Horizonte (idcid=1).
\begin{lstlisting}
UPDATE TIPOQUARTO SET preco = preco * 1.15
WHERE idhot IN (SELECT idhot FROM HOTEL WHERE idcid = 1);
\end{lstlisting}

\clearpage

% ==============================================================================
% ETAPA 6
% ==============================================================================
\section{Etapa 6: Visões (Views)}

\textbf{Visão 1: Hotéis 5 Estrelas (Integridade)}
\begin{lstlisting}
CREATE VIEW view_hoteis_5_estrelas AS
SELECT idhot, idcid, estrelas, enderecohot, nomehot FROM HOTEL
WHERE estrelas = 5 WITH LOCAL CHECK OPTION;
\end{lstlisting}

\textbf{Visão 2: Catálogo Turístico (Disponibilidade)}
\begin{lstlisting}
CREATE VIEW view_catalogo_turistico AS
SELECT PT.idtur, PT.descricaotur, C.nomecid, C.estado
FROM PONTOTURISTICO PT JOIN CIDADE C ON PT.idcid = C.idcid;
\end{lstlisting}

\textbf{Visão 3: Fundadores (Confidencialidade)}
\begin{lstlisting}
CREATE VIEW view_fundadores_publico AS
SELECT nome AS nome_fundador, nacionalidade, atividadepro AS profissao
FROM FUNDADOR;
\end{lstlisting}

\clearpage

% ==============================================================================
% ETAPA 7
% ==============================================================================
\section{Etapa 7: Transações}

\subsection{7(a) Transações Concorrentes}
Simulação de isolamento onde o Terminal 2 não visualiza o \texttt{UPDATE} do Terminal 1 até que ocorra o \texttt{COMMIT}.

\subsection{7(b) Simulação de Deadlock}
Ocorrência de bloqueio mútuo ao tentar atualizar registros cruzados (\textit{Hotel 1} e \textit{Hotel 2}) simultaneamente em duas sessões.

\textbf{Erro gerado pelo PostgreSQL:}
\begin{verbatim}
ERROR:  deadlock detected
Process 215 waits for ShareLock on transaction 779; blocked by process 91.
Process 91 waits for ShareLock on transaction 780; blocked by process 215. 

SQL state: 40P01
Detail: Process 215 waits for ShareLock on transaction 779; blocked by process 91.
Process 91 waits for ShareLock on transaction 780; blocked by process 215.
Hint: See server log for query details.
Context: while updating tuple (0,1) in relation "hotel"
\end{verbatim}

\end{document}